\section{Limit (極限)}
\ssc{Convention}
\sssc{Extended real number system}
When using the extended real number system as the codomain of limits:
\[\forall a\in\bbR\land a>0\colon a/0=\infty,\]
\[\forall a\in\bbR\land a<0\colon a/0=-\infty,\]
when the numerator $0$ is a limit and the result is to be put in a limit or assigned as a result of a limit by limit definition.

When using the extended real number system as the codomain of limits, any function $f$ with limit at $a\in\{-\infty,\infty\}$ being $\pm\infty$ is extended by $f(a)=\pm\infty$ when the input $a$ is a limit and the result is to be put in a limit or assigned as a result of a limit by limit definition.
\sssc{Notation and reading}
$\lim_{x\to a}f(x)=L$ is also denoted as $f(x)\to L$ as $x\to a$, and read $f(x)$ approaches $L$ (for $L=\infty$ also $f(x)$ gets arbitrarily large and for $L=-\infty$ also $f(x)$ gets arbitrarily small) as $x$ approaches (or gets arbitrarily close to) $a$ (for $a=\infty$ also $x$ gets arbitrarily large and for $a=-\infty$ also $x$ gets arbitrarily small).
\subsection{Of Real Sequence}
\subsubsection{Limit}
For a real sequence \(\langle a_n\rangle_{n=l}^{\infty}\), the limit of it (as $n$ approaches infinity) is denoted as $\lim_{n \to \infty} a_n$ or $\lim_n a_n$. For $L\in\bbR$:
\[\lim_{n \to \infty} a_n = L \equiv \forall \varepsilon > 0:\, \exists M \in\mathbb{N}\text{\ s.t.\ } n\in\mathbb{N}\land n \geq M\implies |a_n - L| < \varepsilon.\]
In other words, as \(n\) becomes arbitrarily large, \(a_n\) gets arbitrarily close to \(L\).

If such $M$ exists for some $L\in\bbR$, we say the limit exists and the sequence converges (收斂) to $L$; otherwise, we say the limit doesn't exist and the sequence diverges (發散).
\[\lim_{n\to \infty}a_n=\infty \equiv \forall M > 0, \exists M \in\mathbb{N} \text{\ s.t.\ } n\in\mathbb{N}\land n \geq M \implies a_n > M.\]
\[\lim_{n\to \infty}a_n=-\infty \equiv \forall M < 0, \exists M \in\mathbb{N} \text{\ s.t.\ } n\in\mathbb{N}\land n \geq M \implies a_n < M.\]
\subsubsection{Limit inferior, infimum limit, limit infimum, liminf, inferior limit, lower limit, or inner limit, and limit superior, supremum limit, limit supremum, limsup, superior limit, upper limit, or outer limit}
For a real sequence \(\langle a_n\rangle_{n=l}^{\infty}\), the limit inferior is defined as
\[\liminf_{a\to\infty}a_n\coloneq\sup_{n\geq l}\inf_{m\geq n}a_m,\]
and the limit superior is defined as
\[\limsup_{a\to\infty}a_n\coloneq\inf_{n\geq l}\sup_{m\geq n}a_m.\]
\subsection{Of Real Series}
\subsubsection{Limit}
For a real series
\[S_n = \sum_{i=l}^n a_i,\]
where \(a_i\) are terms of a real sequence \(\langle a_n\rangle_{n=l}^{\infty}\). The limit of \(S_n\) is denoted as \(\lim_{n\to\infty}S_n\) or \(\sum_{i=1}^{\infty}a_i\) and defined as the limit of the sequence \(\left\langle S_n\right\rangle_{n=l}^{\infty}\). If the limit exists and equal to $L$, we say the series converges to $L$; otherwise, we say the series diverges.
\subsubsection{Absolute convergence and conditional convergence}
We say a real series $S_n=\sum_{i=1}^{\infty}a_i$ converges absolutely to $L$ iff the series $\sum_{i=1}^{\infty}\abs{a_i}$ converges to $L$. If a real series is convergent but not convergent absolutely, we say it is convergent conditionally.
\subsection{Of real-valued net}
Below, we will discuss limits of nets in the set of real number.
\subsubsection{Limit}
For a real-valued net $\langle x_a\rangle_{a\in A}$ in which $A$ is a directed set, the limit of $\langle x_a\rangle_{a\in A}$ is denoted as $\lim_a x_a$ or $\lim_{a\to\infty} x_a$. For $L\in\bbR$:
\[\lim_ax_a = L \equiv \forall \varepsilon > 0:\, \exists a_0 \in A\text{\ s.t.\ } a\in A\land a\ge a_0\implies |x_a - L| < \varepsilon.\]

If such $a_0$ exists for some $L\in\bbR$, we say the limit exists or the net converges to $L$; otherwise, we say the limit doesn't exist.
\[\lim_{a}x_a=\infty \equiv \forall M > 0, \exists a_0 \in A \text{\ s.t.\ } a\in A\land a\geq a_0 \implies x_a > M.\]
\[\lim_{a}x_a=-\infty \equiv \forall M < 0, \exists a_0 \in A \text{\ s.t.\ } a\in A\land a\geq a_0 \implies x_a < M.\]
\subsubsection{Limit inferior, infimum limit, limit infimum, liminf, inferior limit, lower limit, or inner limit, and limit superior, supremum limit, limit supremum, limsup, superior limit, upper limit, or outer limit}
For a real-valued net $\langle x_a\rangle_{a\in A}$, the limit inferior is defined as
\[\liminf_{x}x_n\coloneq\sup_{n\geq l}\inf_{m\geq n}x_m,\]
and the limit superior is defined as
\[\limsup_{x}x_n\coloneq\inf_{n\geq l}\sup_{m\geq n}x_m.\]
\subsection{Of Real-Domain Function}
\subsubsection{Limit at a point}
Let \(I\) be an interval containing the point \(a\). Let \( f(x) \) be a function defined on \(I\), except possibly at \(a\) itself. The limit of \( f(x) \) as \( x \) approaches \( a \) is denoted as $\lim_{x \to a} f(x)$. For $L\in\bbR$:
\[\lim_{x \to a} f(x) = L \equiv \forall \varepsilon > 0:\, \exists \delta > 0 \text{\ s.t.\ } 0 < |x - a| < \delta \implies |f(x) - L| < \varepsilon.\]
In other words, as \(x\) becomes arbitrarily close to \(a\), \(f(x)\) gets arbitrarily close to \(L\).

If such $\delta$ exist for some $L\in\bbR$, we say the limit exists; otherwise, we say the limit doesn't exist.
\[\lim_{x\to a}f(x)=\infty \equiv \forall N > 0:\, \exists \delta > 0 \text{\ s.t.\ } 0 < |x - a| < \delta \implies f(x) > N.\]
\[\lim_{x\to a}f(x)=-\infty \equiv \forall N < 0:\, \exists \delta > 0 \text{\ s.t.\ } 0 < |x - a| < \delta \implies f(x) < N.\]
\subsubsection{Limit at Infinity}
Let \(I\) be a left-bounded, right-unbounded interval with the point \(a\) being its endpoint on the left. Let \( f(x) \) be a function defined on \(I\). The limit of \( f(x) \) as \( x \) approaches \( \infty \) is denoted as $\lim_{x \to \infty} f(x)$. For $L\in\bbR$:
\[\lim_{x \to \infty} f(x) = L \equiv \forall \varepsilon > 0: \, \exists M > a \text{\ s.t.\ } x > M \implies |f(x) - L| < \varepsilon.\]
In other words, as \(x\) becomes arbitrarily large, \(f(x)\) gets arbitrarily close to \(L\).

If such $M$ exists for some $L\in\bbR$, we say the limit exists and the function converges to $L$ as \(x\) becomes arbitrarily large; otherwise, we say the limit doesn't exist and the function diverges as \(x\) becomes arbitrarily large.

Let \(I\) be a right-bounded, left-unbounded interval with the point \(a\) being its endpoint on the right. Let \( f(x) \) be a function defined on \(I\). The limit of \( f(x) \) as \( x \) approaches \( -\infty \) is defined as follows:
\[\lim_{x \to -\infty} f(x) = L \equiv \forall \varepsilon > 0: \, \exists M < a \text{\ s.t.\ } x < M \implies |f(x) - L| < \varepsilon.\]
In other words, as \(x\) becomes arbitrarily small, \(f(x)\) gets arbitrarily close to \(L\).

If such $M$ exists for some $L\in\bbR$, we say the limit exists and the function converges to $L$ as \(x\) becomes arbitrarily small; otherwise, we say the limit doesn't exist and the function diverges as \(x\) becomes arbitrarily small.
\[\lim_{x\to\infty}f(x)=\infty \equiv \forall N > 0:\, \exists M > 0 \text{\ s.t.\ } x > M \implies f(x) > N.\]
\[\lim_{x\to\infty}f(x)=-\infty \equiv \forall N > 0:\, \exists M > 0 \text{\ s.t.\ } x > M \implies f(x) < N.\]
\[\lim_{x\to-\infty}f(x)=\infty \equiv \forall N > 0:\, \exists M < 0 \text{\ s.t.\ } x < M \implies f(x) > N.\]
\[\lim_{x\to-\infty}f(x)=-\infty \equiv \forall N > 0:\, \exists M < 0 \text{\ s.t.\ } x < M \implies f(x) < N.\]

$\lim_{x\to\infty}f(x)$ is sometimes denoted as $f(\infty)$; $\lim_{x\to-\infty}f(x)$ is sometimes denoted as $f(-\infty)$.
\subsubsection{One-sided Limits}
\tb{Right-hand Limit (右極限)}: Let \(I\) be a left-open interval with the point \(a\) being its endpoint on the left. Let \( f(x) \) be a function defined on \(I\). The right-hand limit of \( f(x) \) as \( x \) approaches \( a \) is denoted as $\lim_{x \to a^+} f(x)$. For $L\in\bbR$:
\[\lim_{x \to a^+} f(x) = L \equiv \forall \varepsilon > 0 :\,\exists \delta > 0 \text{\ s.t.\ } 0 < x - a < \delta \implies |f(x) - L| < \varepsilon.\]
In other words, as \(x\) becomes arbitrarily close to \(a\) and is greater than \(a\), \(f(x)\) gets arbitrarily close to \(L\).
\[\lim_{x\to a^+}f(x)=\infty \equiv \forall N > 0:\, \exists \delta > 0 \text{\ s.t.\ } 0 < x - a < \delta \implies f(x) > N.\]
\[\lim_{x\to a^+}f(x)=-\infty \equiv \forall N < 0:\, \exists \delta > 0 \text{\ s.t.\ } 0 < x - a < \delta \implies f(x) < N.\]
For a function $f(x)$, $\lim_{x\to a^+}f(x)$ can also be denoted as $f(a^+)$.

$\lim_{x\to-\infty^+}$ is the same as $\lim_{x\to-\infty}$.

\tb{Left-hand Limit (左極限)}: Let \(I\) be a right-open interval with the point \(a\) being its endpoint on the right. Let \( f(x) \) be a function defined on \(I\). The left-hand limit of \( f(x) \) as \( x \) approaches \( a \) is denoted as $\lim_{x \to a^-} f(x)$. For $L\in\bbR$:
\[\lim_{x \to a^-} f(x) = L \equiv \forall \varepsilon > 0 :\,\exists \delta > 0 \text{\ s.t.\ } 0 < a - x < \delta \implies |f(x) - L| < \varepsilon.\]
In other words, as \(x\) becomes arbitrarily close to \(a\) and is less than \(a\), \(f(x)\) gets arbitrarily close to \(L\).
\[\lim_{x\to a^-}f(x)=\infty \equiv \forall N > 0:\, \exists \delta > 0 \text{\ s.t.\ } 0 < a - x < \delta \implies f(x) > N.\]
\[\lim_{x\to a^-}f(x)=-\infty \equiv \forall N < 0:\, \exists \delta > 0 \text{\ s.t.\ } 0 < a - x < \delta \implies f(x) < N.\]
For a function $f(x)$, $\lim_{x\to a^-}f(x)$ can also be denoted as $f(a^-)$.

$\lim_{x\to\infty^-}$ is the same as $\lim_{x\to\infty}$.
\subsubsection{Horizontal asymptote (水平漸近線)}
We say $y=L$ is a horizontal asymptote of $y=f(x)$ iff $\lim_{x \to \infty} f(x)=L$ or $\lim_{x \to -\infty} f(x)=L$.
\subsubsection{Slant asymptote (斜漸近線)}
We say $y=mx+b$ with $m\neq 0$ is a slant asymptote of $y=f(x)$ iff $\lim_{x \to \infty} \qty(f(x)-(mx+b))=0$ or $\lim_{x \to -\infty} \qty(f(x)-(mx+b))=0$.
\subsubsection{Vertical asymptote (鉛直漸近線)}
We say $x=a$ is a vertical asymptote of $y=f(x)$ iff $\abs{\lim_{x \to a^+} f(x)}=\infty$ or $\abs{\lim_{x \to a^-} f(x)}=\infty$.
\subsection{Of Real-valued Functions from Topological Space}
\subsubsection{Limit}
Let $(X,\tau)$ be a topological space, $a\in E\subseteq X$, and $f$ be a function definited on $E$ except possibly at $a$ with codomain $\bbR$. The limit of $f$ as $x$ approaches $a\in E$ is denoted as $\lim_{x\to a}f(x)$. For $L\in\bbR$:
\[\lim_{x \to a} f(x) = L \equiv \forall \varepsilon > 0:\, \exists U\in\tau\land a\in U \text{\ s.t.\ } x\in E\cap U\setminus\{a\} \implies |f(x) - L| < \varepsilon.\]
If such $U$ exist for some $L\in\bbR$, we say the limit exists; otherwise, we say the limit doesn't exist.
\[\lim_{x\to a}f(x)=\infty \equiv \forall N > 0:\, \exists U\in\tau\land a\in U \text{\ s.t.\ }x\in E\cap U\setminus\{a\}\implies f(x) > N.\]
\[\lim_{x\to a}f(x)=\infty \equiv \forall N < 0:\, \exists U\in\tau\land a\in U \text{\ s.t.\ }x\in E\cap U\setminus\{a\}\implies f(x) < N.\]
\subsubsection{Limit inferior, infimum limit, limit infimum, liminf, inferior limit, lower limit, or inner limit, and limit superior, supremum limit, limit supremum, limsup, superior limit, upper limit, or outer limit}
Let $(X,\tau)$ be a topological space, $a\in E\subseteq X$, and $f$ be a function definited on $E$ except possibly at $a$ with codomain $\bbR$. The limit inferior of $f$ as $x$ approaches $a\in E$ is defined as
\[\liminf_{x\to a}f(x)\coloneq\sup\{\inf\{f(x)\mid x\in E\cap U\setminus\{a\}\}\mid a\in U\in\tau\},\]
and the limit superior of $f$ as $x$ approaches $a\in E$ is defined as
\[\limsup_{x\to a}f(x)\coloneq\inf\{\sup\{f(x)\mid x\in E\cap U\setminus\{a\}\}\mid a\in U\in\tau\}.\]
\subsection{Of function between normed spaces}
\sssc{Definition}
Let $V,W$ be two normed spaces. Let $I$ be an open subset of $V$ containing the point $a$. Let $f$ be a function with codomain $W$ defined on $I$, except possibly at $a$ itself. The limit of $f(x)$ as $x$ approaches $a$ is denoted as $\lim_{x\to a}f(x)$ or $\lim_{\|x-a\|_V\to0}f(x)$, and for $x=(x_1,x_2,\ldots,x_n)$ and $a=(a_1,a_2,\ldots,a_n)$ also $\lim_{(x_1,x_2,\ldots,x_n)\to(a_1,a_2,\ldots,a_n)}f(x_1,x_2,\ldots,x_n)$, $\lim_{\substack{x_1\to a_1\\x_2\to a_2\\\vdots\\x_n\to a_n}}f(x_1,x_2,\ldots,x_n)$. For $L\in W$:
\[\lim_{x\to a}f(x)=L\equiv\forall\varepsilon\exists\delta>0\tx{\ s.t.\ }0<\|x-a\|_V<\delta\implies\|f(x)-L\|_W<\varepsilon.\]
If such $M$ exist for some $L\in W$, we say the limit exists; otherwise, we say the limit doesn't exist.
\subsection{Of Nets of Fuctions}
\subsubsection{(Pointwise) limit}
Let $X$ be a set and $Y$ be a topological space. A net of functions $\langle f_a\rangle_{a\in A}$ all having the same domain $X$ and codomain $Y$ is said to converge pointwise (逐點收斂) to a function $f\colon X\to Y$ or have a pointwise limit function $f\colon X\to Y$, denoted as 
\[\lim_af_a=f\text{\ pointwise}\]
or
\[\lim_af_a=f,\]
iff
\[\forall x\in X\colon\lim_af_a(x)=f(x).\]
\subsubsection{Uniform limit}
Let $X$ be a set and $(Y,d)$ be a metric space. A net of functions $\langle f_a\rangle_{a\in A}$ all having the same domain $X$ and codomain $Y$ is said to converge uniformly (一致收斂) to a function $f\colon X\to Y$ or have a uniform limit function $f\colon X\to Y$, denoted as 
\[\lim_af_a=f\text{\ uniformly},\]
iff
\[\forall\varepsilon\in\mathbb{R}_{>0}\colon\exists N\in A\colon a\in A\land N\leq a\implies\sup_{x\in X}d(f_a(x),f(x))<\varepsilon.\]
\subsection{Of Nets of Fuctions on Discrete Points to Function on Connected Space}
\subsubsection{(Pointwise) limit}
Let $X$ be a connected topological space and $Y$ be a topological space. A net of functions $\langle f_a\rangle_{a\in A}$ having domains $X_a\subseteq X$ and codomain $Y$ is said to converge pointwise to a function $f\colon X\to Y$ or have a pointwise limit function $f\colon X\to Y$, denoted as 
\[\lim_af_a=f\text{\ pointwise}\]
or
\[\lim_af_a=f,\]
iff
\[\forall x\in X\forall\langle x_a\rangle_{a\in A}\colon\qty(\qty(\forall a\in A\colon x_a\in X_a)\land\lim_ax_a=x)\implies\lim_af_a(x_a)=f(x).\]
\subsubsection{Uniform limit}
Let $X$ be a connected topological space and $(Y,d)$ be a metric space. A net of functions $\langle f_a\rangle_{a\in A}$ having domains $X_a\subseteq X$ and codomain $Y$ is said to converge uniformly to a function $f\colon X\to Y$ or have a uniform limit function $f\colon X\to Y$, denoted as 
\[\lim_af_a=f\text{\ uniformly},\]
iff
\[\forall\varepsilon\in\mathbb{R}_{>0}\exists N\in A\forall b\in A\forall\langle x_a(x)\rangle_{a\in A}\colon\qty(\forall a\in A\colon x_a(x)\in X_a)\land\qty(\forall x\in X\colon\lim_ax_a(x)=x)\land N\leq b)\implies\sup_{x\in X}d(f_b(x_b),f(x))<\varepsilon).\]
\ssc{Big O notation and little notation}
\sssc{Big O notation at a point}
For a point $a$ in a topological space and real-valued functions $f$ and $g$ defined on some neighborhood of $a$ except possibly at $a$, $f(x)=O(g(x))$ as $x\to a$ if there exists some neighborhood $I$ of $a$ and positive real number $M$ such that 
\[\abs{f(x)}\leq M\abs{g(x)}.\]
\sssc{Little o notation at a point}
For a point $a$ in a normed space and real-valued functions $f$ and $g$ defined on some neighborhood of $a$ except possibly at $a$, $f(x)=o(g(x))$ as $x\to a$ if
\[\lim_{x\to a}\frac{f(x)}{g(x)}=0.\]
\sssc{Big O notation at infinity}
For real-valued functions $f$ and $g$ defined on some interval $[a,\infty)$, $f(x)=O(g(x))$ as $x\to\infty$ if there exists some real number $T$ such that for all $x\geq T$,
\[\abs{f(x)}\leq\abs{g(x)}.\]
For real-valued functions $f$ and $g$ defined on some interval $(-\infty,b]$, $f(x)=O(g(x))$ as $x\to-\infty$ if there exists some real number $T$ such that for all $x\leq T$,
\[\abs{f(x)}\leq\abs{g(x)}.\]
\sssc{Little o notation at infinity}
For real-valued functions $f$ and $g$ defined on some interval $[a,\infty)$, $f(x)=o(g(x))$ as $x\to\infty$ if
\[\lim_{x\to\infty}\frac{f(x)}{g(x)}=0.\]
For real-valued functions $f$ and $g$ defined on some interval $(-\infty,b]$, $f(x)=o(g(x))$ as $x\to-\infty$ if
\[\lim_{x\to-\infty}\frac{f(x)}{g(x)}=0.\]